\documentclass{article}
\usepackage{hyperref,amsmath}
\begin{document}
\section*{When a reconstructed reply is sufficient to resume a retry}
\paragraph{Scope.} This is a source-level comparison for a digital old
duty indexed by a stable request and lease, not a finding about a
running server. The intentionally premature execution order is
outside RIFL's documented rendezvous premise. The fixture checks
why a simulator adapter must not present that order as a supported
recovery path. Physical clearance, durable replication, authentication,
bounded network timing and the mechanism that establishes cut closure
are outside this test.

\paragraph{Source boundary.}
On the pinned RAMCloud revision from PR 853, the extracted
\texttt{shouldRecover}, \texttt{recoverRecord},
\texttt{checkDuplicate} and \texttt{processAck} methods are exercised
against isolated serial state. Restoring a qualified old completion
before lookup returns the exact saved result. If a fixture performs
ordinary new-request classification on that key \emph{before}
restoring it, it installs an in-progress entry; the subsequent import
does not turn that entry into a saved completion. The same import
after an independent lane's request succeeds normally and preserves
both exact results. This test does not execute an old mutation in
the premature branch. It distinguishes readiness of the named
rendezvous from readiness of the process, and applies the established
requirement that completion records precede retry handling
(\href{https://doi.org/10.1145/2815400.2815416}{Lee et al., RIFL,
Sections 3--4}). A whole-station pause is not the strongest
baseline when disjoint readiness scopes are available.

\paragraph{Matched possibility boundary for incomplete reconstruction.}
Assume a single fenced cut, a stable key, at most one lawful immutable
result, and a finite exhaustive family $\mathcal C$ of \emph{nonempty
sets of complete copy holders} for histories where that result exists.
Atomic effect and record, plus preservation of at least one complete
copy in every admitted committed history, are explicit premises. If a
committed effect can leave no readable record anywhere, the empty
transcript cannot distinguish it even after checking every holder.
All these histories can otherwise yield the same available evidence
as one admitted empty history. On a subset $S$ of holders, completed
negative reads govern the same cut and correctly report the absence of
the keyed record there. Then the all-negative transcript rules out
prior completion exactly if $S$ meets every member of $\mathcal C$.
Sufficiency follows by intersection. For necessity choose an
admitted holder set disjoint from $S$; its transcript and that of
the empty history agree, but safe initiation of an otherwise
authorized non-idempotent effect differs. A positive read of an
authentic immutable result can answer the exact old query without
waiting for unresolved negative reads. It grants no authority to
perform a new effect, nor does a proven empty result itself grant
that authority.

The worst-case number of unit-cost reads for the binary existence
decision, with all holders readable and no further information
before a response, is the minimum hitting-set size. A fixed
minimum hitting set attains it; the run of an arbitrary adaptive
algorithm against the empty history forces its read set to hit
every member. Thus if any pair of three holders is possible, two
qualified negatives suffice. If a changed recovery mode can leave
the exact old result on any singleton, the old two-read negative
no longer suffices and three are necessary. If a holder cannot be
read and no accessible hitting set survives, an all-negative
transcript cannot license a new effect. It can still leave other
object scopes operational. Two negative reads at different
unfenced cuts also fail: check A while the only copy is at B,
then check B after the copy moves to A. A same-cut premise is
essential even if each local reply was truthful.

\paragraph{Comparison and checks.}
The intersection rule is ordinary read/write quorum reasoning
(\href{https://doi.org/10.1145/800215.806583}{Gifford 1979}).
\href{https://lamport.org/pubs/reconfiguration-tutorial.pdf}{Lamport,
Malkhi and Zhou} provide established state-machine reconfiguration
comparators where replication and configuration transitions are
implemented. The relevant R2 obligation here is to re-establish
which copy sets can serve a named exact old response after a
changed recovery dependency, rather than to transplant the
old negative certificate by convention. The simplified scope is
checked on 51 matched source controls and all 524,272 negative
projections of nonempty families on four possible holders.
The copy laws used for the mathematical examples are stipulated;
the RAMCloud fixture does not establish them or improve on an
equally informed quorum and rendezvous baseline.
\end{document}
